\chapter{Manual de usuario}\label{manual}

\section{Introducci\'on}

Este cap\'itulo describe el uso operativo del sistema desde la perspectiva de sus
usuarios finales. El objetivo no es detallar la implementaci\'on interna,
sino proporcionar una gu\'ia clara para ejecutar las tareas m\'as frecuentes
de forma correcta y segura.

El sistema dispone de tres perfiles principales:
\textbf{administrador}, \textbf{profesor} y \textbf{estudiante}.
Cada perfil accede a vistas y operaciones diferentes seg\'un permisos.

\section{Acceso al sistema}

\subsection{Requisitos previos de uso}

Antes de operar con normalidad en la plataforma,
se recomienda verificar las siguientes condiciones:

\begin{itemize}
	\item disponer de navegador actualizado en versi\'on estable;
	\item comprobar conectividad de red durante subida de archivos de gran tama\~no;
	\item utilizar credenciales personales y no compartidas;
	\item validar que el perfil asignado (administrador/profesor/estudiante) coincide con el rol esperado.
\end{itemize}

Estas precondiciones reducen incidencias de sesi\'on,
errores de permisos y fallos de subida de materiales.

\subsection{Pantallas de acceso}

El acceso se realiza desde la pantalla de inicio de sesi\'on.
Las rutas funcionales principales son:

\begin{itemize}
	\item \texttt{/login}: inicio de sesi\'on.
	\item \texttt{/forgot-password}: solicitud de recuperaci\'on de contrase\~na.
	\item \texttt{/reset-password}: establecimiento de nueva contrase\~na mediante token.
\end{itemize}

Tras autenticarse correctamente, el usuario es redirigido al panel principal
(\texttt{/dashboard}). Si el usuario ya tiene sesi\'on activa, la ruta de inicio
redirige autom\'aticamente al panel.

\subsection{Ciclo de sesi\'on}

La sesi\'on se basa en token de acceso y token de refresco.
Cuando el token de acceso expira, el cliente intenta renovarlo autom\'aticamente.
Si la renovaci\'on falla, la sesi\'on se cierra y el usuario debe autenticarse de nuevo.

\subsection{Recuperaci\'on de contrase\~na}

El flujo de recuperaci\'on es el siguiente:

\begin{enumerate}
	\item El usuario solicita recuperaci\'on indicando su correo.
	\item El sistema emite un token temporal y env\'ia instrucciones de restablecimiento.
	\item El usuario accede al formulario de nueva contrase\~na e introduce token y clave nueva.
\end{enumerate}

\section{Perfiles y permisos}

La Tabla~\ref{tab:manual-perfiles} resume capacidades por perfil.

\begin{table}[ht]
\centering
\begin{tabular}{|P{2.6cm}|P{9.3cm}|}
\hline
\textbf{Perfil} & \textbf{Operaciones principales} \\
\hline
Administrador & Gestionar usuarios, cursos y grupos; asignar o retirar roles; mantener estructura acad\'emica. \\
\hline
Profesor & Crear/editar actividades; crear/editar entregables; revisar entregas; calificar; publicar/ocultar notas. \\
\hline
Estudiante & Consultar cursos y actividades visibles; realizar entregas; consultar historial, feedback y calificaciones seg\'un visibilidad. \\
\hline
\end{tabular}
\caption{Capacidades por perfil de usuario}
\label{tab:manual-perfiles}
\end{table}

\section{Gu\'ia de uso para profesorado}

\subsection{Crear y configurar una actividad}

Flujo recomendado:

\begin{enumerate}
	\item Acceder a \texttt{/cursos} y seleccionar un curso.
	\item Crear actividad desde el detalle del curso.
	\item Definir t\'itulo, descripci\'on, fechas y visibilidad.
	\item Asignar grupos destinatarios cuando aplique.
\end{enumerate}

Validaciones clave:

\begin{itemize}
	\item El t\'itulo es obligatorio.
	\item La fecha de inicio no puede ser posterior a la fecha l\'imite.
	\item La visibilidad condiciona si el alumnado puede ver la actividad.
\end{itemize}

\subsection{Crear y editar entregables}

Desde la actividad, el profesor puede crear entregables
desde la vista de alta de entregables del detalle de actividad.

Par\'ametros relevantes de configuraci\'on:

\begin{itemize}
	\item T\'itulo, descripci\'on y fecha l\'imite.
	\item Tipo de archivo esperado (PDF, DOCX, ZIP, imagen, video, texto, enlace, etc.).
	\item Tama\~no m\'aximo permitido.
	\item Permitir o no reenv\'io.
	\item Visibilidad de notas para estudiantes.
\end{itemize}

Cuando el tipo esperado es ZIP, puede definirse estructura esperada
y validaci\'on estricta, lo que obliga a que la entrega siga el formato indicado.

\subsection{Revisar entregas y evaluar}

Las entregas pueden revisarse desde:

\begin{itemize}
	\item \texttt{/entregables/:id} (listado por entregable),
	\item \texttt{/evaluaciones} (bandeja de pendientes de calificaci\'on).
\end{itemize}

Acciones habituales:

\begin{enumerate}
	\item Filtrar por curso, actividad, entregable y estado temporal.
	\item Abrir detalle de entrega (\texttt{/entregas/:id}).
	\item Revisar archivos, comentario y metadatos.
	\item Registrar nota y feedback.
	\item Publicar/ocultar visibilidad de nota al estudiante.
\end{enumerate}

El sistema permite descargar entregas de forma agregada en ZIP,
tanto a nivel de entregable como de actividad completa.

\subsection{Flujo docente recomendado de extremo a extremo}

Para maximizar consistencia y reducir retrabajo,
se propone la siguiente secuencia de trabajo docente:

\begin{enumerate}
	\item crear actividad con fechas y grupos correctamente delimitados;
	\item publicar entregables con requisitos de formato y tama\~no expl\'icitos;
	\item revisar peri\'odicamente estado de entregas pendientes;
	\item corregir por lotes homog\'eneos (por grupo o por entregable);
	\item registrar feedback accionable y criterio de nota transparente;
	\item publicar calificaciones cuando el bloque de correcci\'on est\'e validado.
\end{enumerate}

Este flujo mejora la trazabilidad y evita inconsistencias
entre estado de evaluaci\'on y comunicaci\'on al alumnado.

\section{Gu\'ia de uso para estudiantes}

\subsection{Consultar actividades y entregables}

Flujo recomendado:

\begin{enumerate}
	\item Acceder al panel y entrar en \texttt{/cursos}.
	\item Seleccionar curso y revisar actividades visibles.
	\item Entrar en el detalle de actividad y localizar entregables disponibles.
\end{enumerate}

La visibilidad de actividad/entregable depende de la configuraci\'on docente
y de la pertenencia del estudiante a grupos autorizados.

\subsection{Realizar una entrega}

La subida se realiza en \texttt{/entregables/:id/entregar}.
El formulario adapta su comportamiento al tipo de entrega:

\begin{itemize}
	\item \textbf{SOLO\_TEXTO}: no admite archivos; requiere comentario.
	\item \textbf{ENLACE}: no admite archivos; el contenido debe ir en comentario.
	\item \textbf{ZIP}: exige archivo ZIP y puede validar estructura.
	\item \textbf{Otros tipos}: valida extensi\'on y tama\~no m\'aximo seg\'un configuraci\'on.
\end{itemize}

Adem\'as, el formulario incluye guardado temporal en local
(borrador), lo que reduce riesgo de p\'erdida de trabajo
si se interrumpe la sesi\'on o falla la conexi\'on.

\subsection{Recomendaciones para entregas sin incidencias}

Para minimizar errores de validaci\'on,
se recomienda al estudiante:

\begin{itemize}
	\item revisar el tipo de entrega requerido antes de preparar el material;
	\item verificar nombre y extensi\'on de archivos antes de subir;
	\item evitar subir archivos en el \`ultimo minuto para mitigar riesgos de red;
	\item confirmar en la vista de detalle que la versi\'on final ha quedado registrada.
\end{itemize}

En entregas de tipo ZIP,
si existe estructura obligatoria,
conviene validar localmente el contenido antes de la subida.

\subsection{Consultar historial y feedback}

En el detalle de entrega (\texttt{/entregas/:id}) el estudiante puede:

\begin{itemize}
	\item revisar versiones entregadas,
	\item descargar/previsualizar materiales,
	\item consultar feedback y nota cuando est\'en publicadas.
\end{itemize}

Si la nota no es visible, la entrega aparece sin calificaci\'on p\'ublica,
aunque internamente ya haya sido evaluada por el profesor.

\section{Gu\'ia de uso para administraci\'on}

\subsection{Panel de administraci\'on}

El panel \texttt{/admin} est\'a organizado en tres pesta\~nas principales:

\begin{itemize}
	\item \textbf{Usuarios}: alta, edici\'on, eliminaci\'on y asignaci\'on de roles.
	\item \textbf{Cursos}: creaci\'on y mantenimiento de cursos.
	\item \textbf{Grupos}: creaci\'on y gesti\'on de grupos asociados a cursos.
\end{itemize}

\subsection{Operaciones destacadas}

El perfil administrador puede:

\begin{enumerate}
	\item Crear usuarios individualmente o por importaci\'on de archivo.
	\item Asignar rol de profesor o estudiante (con grupo en el caso de estudiante).
	\item Crear cursos y vincular profesorado.
	\item Estructurar grupos y reasignar alumnado.
\end{enumerate}

Este perfil concentra operaciones de gobierno funcional,
por lo que se recomienda limitar su uso a personal autorizado.

\subsection{Pol\'iticas recomendadas de gobierno de usuarios}

Para mantener orden operativo y seguridad,
se recomienda aplicar las siguientes pol\'iticas de administraci\'on:

\begin{itemize}
	\item principio de m\'inimo privilegio en asignaci\'on de roles;
	\item revisi\'on peri\'odica de cuentas inactivas o duplicadas;
	\item separaci\'on de cuentas de administraci\'on y cuentas de uso docente;
	\item registro de cambios sensibles en altas, bajas y cambios de rol.
\end{itemize}

\section{Mensajes de error y resoluci\'on de incidencias}

Las incidencias m\'as frecuentes y su acci\'on recomendada son:

\begin{itemize}
	\item \textbf{Credenciales inv\'alidas}: verificar correo y contrase\~na.
	\item \textbf{Sesi\'on expirada}: iniciar sesi\'on de nuevo.
	\item \textbf{Archivo no permitido}: revisar tipo/extensi\'on/tama\~no del entregable.
	\item \textbf{Sin permisos}: comprobar rol y pertenencia a curso/grupo.
	\item \textbf{Error al enviar entrega}: reintentar y comprobar conexi\'on.
\end{itemize}

Como recomendaci\'on general, ante errores persistentes debe registrarse
la hora, operaci\'on realizada y mensaje mostrado para facilitar diagn\'ostico.

\subsection{Matriz de incidencias operativas}

\begin{table}[ht]
\centering
\begin{tabular}{|P{3.2cm}|P{4.4cm}|P{3.8cm}|}
\hline
\textbf{Incidencia} & \textbf{Causa habitual} & \textbf{Acci\'on recomendada} \\
\hline
No se puede iniciar sesi\'on & Credenciales incorrectas o sesi\'on expirada & Reintentar autenticaci\'on y, si procede, usar recuperaci\'on de contrase\~na \\
\hline
Entrega rechazada & Tipo de archivo o tama\~no no permitido & Revisar configuraci\'on del entregable y ajustar archivo \\
\hline
No aparece actividad/entregable & Visibilidad desactivada o grupo no asignado & Verificar configuraci\'on por parte del profesorado \\
\hline
No se puede calificar & Falta de permisos por rol/contexto & Confirmar vinculaci\'on docente con curso/grupo \\
\hline
Error en descarga ZIP & Interrupci\'on de red o contenido no disponible & Reintentar descarga y revisar estado de materiales \\
\hline
\end{tabular}
\caption{Matriz de incidencias frecuentes y resoluci\'on recomendada}
\label{tab:manual-incidencias}
\end{table}

\section{Seguridad y protecci\'on de la informaci\'on en el uso diario}

La seguridad operativa depende tanto de mecanismos t\'ecnicos
como de buenas pr\'acticas de uso.
Para preservar confidencialidad e integridad de datos acad\'emicos,
se recomienda:

\begin{itemize}
	\item no compartir contrase\~nas ni tokens de sesi\'on;
	\item cerrar sesi\'on en dispositivos compartidos;
	\item evitar cargas de archivos desde equipos no confiables;
	\item revisar peri\'odicamente permisos y visibilidad de notas.
\end{itemize}

Estas medidas reducen riesgos de acceso no autorizado
y publicaci\'on accidental de informaci\'on acad\'emica sensible.

\section{Checklist operativo por perfil}

Como apoyo de uso r\'apido,
se incluye un checklist resumido para cada rol:

\begin{itemize}
	\item \textbf{Administrador}: alta de usuarios, asignaci\'on de roles, revisi\'on de estructura de cursos/grupos.
	\item \textbf{Profesor}: configuraci\'on de actividad, publicaci\'on de entregable, correcci\'on y publicaci\'on de nota.
	\item \textbf{Estudiante}: revisi\'on de actividad visible, entrega en plazo, verificaci\'on de estado y feedback.
\end{itemize}

El uso sistem\'atico de este checklist disminuye errores de operaci\'on
y mejora la consistencia del proceso acad\'emico.

\section{Buenas pr\'acticas de uso}

Para minimizar incidencias operativas:

\begin{itemize}
	\item comprobar siempre fecha l\'imite y tipo de entrega antes de enviar,
	\item revisar la vista de detalle tras cada acci\'on cr\'itica,
	\item mantener actualizado el perfil de usuario,
	\item evitar compartir credenciales,
	\item cerrar sesi\'on en equipos compartidos.
\end{itemize}

\section{Conclusiones del cap\'itulo}

El manual de usuario consolida el uso operativo del sistema y garantiza
que los tres perfiles principales disponen de un procedimiento claro,
coherente con las reglas de negocio y alineado con los requisitos definidos.

Su existencia facilita adopci\'on, reduce errores de uso y mejora la calidad
del servicio acad\'emico prestado por la plataforma.

